\documentclass[11pt,a4paper]{article}

\usepackage[T1]{fontenc}
\usepackage[utf8]{inputenc}
\usepackage[margin=1in]{geometry}
\usepackage{amsmath}
\usepackage{booktabs}
\usepackage{natbib}
\usepackage{hyperref}
\hypersetup{
  colorlinks = true,
  linkcolor  = black,
  citecolor  = black,
  urlcolor   = blue
}
\usepackage{orcidlink}

\title{\texttt{attest}: Fail-Closed Machine Learning with Model
Certificates in \textsf{R}}

\author{%
  Diogo Ribeiro\,\orcidlink{0009-0001-2022-7072}\\[2pt]
  \small Faculty of Media Arts and Design, Technical University of Porto,
  Portugal\\
  \small \texttt{dfr@esmad.ipp.pt}
}

\date{22 September 2026}

\begin{document}
\maketitle

\begin{abstract}
A fitted model in \textsf{R} will answer any question put to it, and the
validation that justified it lives somewhere else entirely. \texttt{attest}
binds the two together: models carry an immutable certificate of the checks
run at fit time, predictions return a validity status alongside a value, and
inputs outside the certified support are refused rather than answered. The
package treats every threshold as a comparison against a null rather than a
fixed line, which materially changes what its checks report: the conventional
population stability index threshold flags roughly 80\% of healthy 50-row
batches, where a simulated null flags 6\%. Where covariates have shifted but
the outcome relationship has not, reweighted conformal intervals recover
coverage from 0.74 to 0.95 against a 0.90 target.
\end{abstract}

\section{Summary}

A fitted model in \textsf{R} will answer any question put to it. Hand
\texttt{predict()} a row from a population the model has never seen, or one
that has drifted for six months, and a number comes back carrying no
indication that it should not be trusted. The validation that justified the
model lives elsewhere---in a script, a notebook, a slide---and nothing binds
it to the object. By the time a model reaches production, the evidence and the
artefact have parted company.

\texttt{attest} closes that gap by making the evidence part of the object and
refusal part of the prediction path. A model fitted with \texttt{attest\_fit()}
carries a \emph{certificate}: leakage, calibration, conformal coverage and
distribution-shift checks that run at fit time and are stored inside the model,
bound to the training data, the specification and the fitted object by SHA-256
hash. Predictions are not bare numbers. Each row returns an interval
(regression) or a conformal label set (classification), a validity status, and
a reason when that status is not \texttt{valid}. Rows falling outside the
certified support are refused rather than answered, and overriding that refusal
is recorded on the returned object. Model cards are rendered from the
certificate, in Markdown or as a self-contained HTML page with diagnostic
charts, so a report cannot disagree with the model it describes.

The package wraps existing modelling tools rather than replacing them.
\texttt{engine\_parsnip()} accepts any \texttt{parsnip} \citep{parsnip}
specification, which reaches \texttt{ranger} \citep{ranger}, \texttt{xgboost}
and \texttt{glmnet} without an adapter for each, and the certificate records
which model was actually fitted.

\section{Statement of need}

Tooling for trustworthy machine learning in \textsf{R} is plentiful but
optional. Calibration, leakage screening and conformal prediction are each
available as separate packages, and each must be remembered, invoked and
interpreted by the analyst. Nothing prevents a model from being deployed
without them, and nothing records afterwards whether they were run. The failure
mode is not that the checks are unavailable; it is that they are easy to skip
and impossible to audit after the fact.

\texttt{attest} inverts the default. Checks execute inside
\texttt{attest\_fit()}, so there is no code path that fits a model without
them. A blocking failure refuses to issue a certificate at all, and a model
whose certificate is not valid cannot predict. Waiving a check is possible but
requires a written reason, which is stored in the certificate and reprinted in
every report thereafter. The package is aimed at regulated and safety-adjacent
settings---credit, insurance pricing, clinical triage---where the question
asked of a model is not only how accurate it is, but what was checked, on what
data, and by whom.

Two design commitments distinguish it from a checklist.

\subsection{Every threshold is compared against a null}

Checks report a bootstrap confidence interval alongside each statistic and fail
only when the whole interval clears the threshold, so sampling noise alone
cannot fail a model; an interval straddling the threshold is reported as
\texttt{weak} rather than silently passing.

This matters in practice. The population stability index, the conventional
shift statistic, is biased upward in small batches: with ten bins and 50 rows
drawn from the training distribution \emph{itself}, its median value is
approximately 0.2, which is the conventional flag threshold. Comparing against
a simulated null at the observed batch size rather than a fixed line reduces
the measured false-positive rate on unshifted 50-row batches---the package
default minimum---from roughly 80\% to 6\%, while leaving detection of a
genuine shift unchanged (Table~\ref{tab:psi}).

\begin{table}[ht]
\centering
\caption{False-positive rate on unshifted batches, three features, batch
flagged if any feature fires.}
\label{tab:psi}
\begin{tabular}{lrr}
\toprule
Batch size & Fixed threshold 0.2 & Simulated null \\
\midrule
50   & 80\% & 6\% \\
60   & 69\% & 6\% \\
200  & 0\%  & 0\% \\
\bottomrule
\end{tabular}
\end{table}

A classifier two-sample test \citep{lopezpaz2017} complements it, detecting
changes in the dependence between features that any per-feature statistic
misses by construction. Its scores are cross-fitted, because an in-sample area
under the curve reads approximately 0.53 on two samples drawn from the same
distribution and would report shift where there is none.

\subsection{Refusal is where the guarantee ends}

Split conformal prediction \citep{vovk2005,lei2018} provides marginal coverage
under exchangeability, and that assumption is void---not merely
weakened---under distribution shift. \texttt{attest} refuses rows outside the
certified support because outside it there is no guarantee left to offer;
refusal is an admission rather than an added safeguard.

Where the covariate distribution has moved but the conditional distribution of
the outcome has not, \texttt{conformal\_split(weighted = TRUE)} reweights the
calibration scores by an estimated density ratio \citep{tibshirani2019},
widening intervals instead of only warning. On a heteroscedastic problem
shifted by 1.5 standard deviations, a fixed quantile covers approximately 0.74
against a 0.90 target while reporting its usual width; the reweighted intervals
cover 0.95. When every weight is one, the weighted quantile reduces exactly to
the ordinary one, so enabling it costs nothing when nothing has moved.

\section{Limitations}

The documentation is explicit about what a certificate does not establish.
Conformal coverage is marginal rather than conditional. It is achieved by
construction for any underlying model and therefore cannot distinguish a
leaking model from an honest one: in a worked example, a model given a feature
that restates the outcome is better calibrated and sharper than an honest one
while its coverage is indistinguishable. Reweighting corrects covariate shift
only and is actively misleading under concept drift, because the intervals look
adjusted. No shift check inspects the outcome, so none of them can distinguish
the input distribution having moved from the model having become worse.

\section*{Acknowledgements}

\texttt{attest} builds on \texttt{cli}, \texttt{digest}, \texttt{jsonlite},
\texttt{rlang} and \texttt{tibble}.

\bibliographystyle{plainnat}
\bibliography{paper}

\end{document}
